\section{Engine study}
\label{sec:engine}

Section~\ref{sec:serving} measured a research serving path without a production engine. Here
we measure \expone{} under vLLM 0.30.0 \citep{kwon2023vllm} on one NVIDIA H100 80GB GPU, the
GPU we chose for serving (our internal deployment runs the same artifact on one L40S 48GB,
because no H100 was available; Section~\ref{sec:availability}). All numbers in this section
were measured on 2026-09-25.

\subsection{Setup}
\label{sec:engine_setup}

\paragraph{Hardware and engine.}
One H100 80GB HBM3 GPU of an eight-GPU node, vLLM 0.30.0 with
Torch 2.13 and CUDA 13. The engine receives the token IDs produced by the Bobcat compiler
unchanged.

\paragraph{Readout.}
We read only the raw logits of the offered identifier tokens (vLLM's raw-logit log-probability
mode with an explicit token list). Every request asks for a fixed width of 64 identifiers; a
question with more candidates is asked several times with the same prompt, and the extra
requests reuse the cached prompt. \textbf{vLLM computes one token per request at temperature
0 and discards it}; that token never enters a reply. This differs from the research path of
Section~\ref{sec:serving}, which samples nothing, so ``no generation'' on this engine means
that no generated token reaches the caller, not that the engine performs no decode step.

\paragraph{Served artifact.}
The LoRA adapter is merged into the base weights with the sum computed in float32 and stored
in BF16. For FP8, vLLM quantizes the merged BF16 weights dynamically at load time. The merged
FP8 model is the artifact our internal deployment serves.

\paragraph{Workloads.}
The workload profiles W1 to W4 of our specification. W1 to W3 are measured warm with 20 repetitions (10 for the 8K and
16K inputs), each repetition on a different document; W4 submits between 1 and 512 distinct,
independent 1,024-token requests. Timings are taken at the engine call and exclude HTTP and the
network.

\paragraph{Engine constraints we hit.}
\begin{itemize}
  \item vLLM's default of 1,024 concurrent sequences exceeds the number of cache blocks it can
    allocate on an 80\,GB H100 for the hybrid layers' recurrent state (which vLLM calls the
    Mamba cache): 292 in BF16 and 763 in FP8. We capped concurrency at 256.
  \item When the number of requested log-probabilities varied between requests, Triton
    compiled a kernel for each new shape and the engine stalled for minutes (observed with
    \texttt{py-spy}). A fixed width of 64 removes this.
  \item The full-vocabulary log-probability path that we first used for questions with more
    than 128 candidates ran out of memory. Chunks of 64 identifiers replace it.
\end{itemize}

\subsection{Quality of the served artifact}
\label{sec:engine_quality}

\begin{table}[t]
\centering
\small
\caption{Development accuracy (\%) of the served artifacts against the evaluated model
(3,188 decisions). Agreement is the fraction of development answers equal to the evaluation
path of the sealed final; gaps are the largest and mean absolute probability differences.}
\label{tab:engine_quality}
\setlength{\tabcolsep}{3.5pt}
\begin{adjustbox}{max width=\linewidth}
\begin{tabular}{@{}lrrrrrrrrr@{}}
\toprule
Artifact & Macro & Search & Citation & Tool & Doc.\ screen & Routing & Class. & Agreement & Max / mean gap \\
\midrule
Evaluated: unmerged LoRA, transformers & 93.4 & 92.2 & 99.8 & 95.0 & 96.5 & 99.3 & 77.6 & -- & -- \\
Merged BF16, vLLM & 93.3 & 92.9 & 100.0 & 95.0 & 96.0 & 98.9 & 76.8 & 98.6\% & 0.90 / 0.013 \\
\textbf{Merged FP8, vLLM (served)} & \textbf{93.1} & 92.6 & 99.8 & 94.5 & 96.0 & 99.1 & 76.4 & \textbf{98.3\%} & 0.80 / 0.015 \\
\bottomrule
\end{tabular}
\end{adjustbox}
\end{table}

The served FP8 artifact scores 93.1\% on development against 93.4\% for the evaluated model,
and 1.7\% of its answers differ (1.4\% for merged BF16; Table~\ref{tab:engine_quality}). The
largest single probability differences are large (0.80 and 0.90) even though the mean
differences are near 0.01. The sealed final result of 94.6\% was measured on the unmerged
path and does not transfer to the served artifact; a separately released served model would
need a new sealed split, because the current final has been opened. The float32 merge did not
remove the merge difference seen in Section~\ref{sec:determinism}; because the merge
procedure and the engine both changed, the two measurements are not a controlled comparison
and we do not attribute the remaining difference.

\subsection{Latency and throughput}
\label{sec:engine_latency}

\begin{table}[t]
\centering
\small
\caption{Latency (p50 in ms unless marked) and throughput on one H100 with vLLM 0.30.0.
Qwen3.5-9B is untrained and measured for speed only. Engine-call timings, no HTTP or network.
W4 reports the best throughput over batch sizes 1 to 512.}
\label{tab:engine}
\setlength{\tabcolsep}{4pt}
\begin{adjustbox}{max width=\linewidth}
\begin{tabular}{@{}lrrrr@{}}
\toprule
Workload & 27B FP8 (served) & 27B BF16 & 9B FP8 (speed only) & 9B BF16 (speed only) \\
\midrule
W1: 512 tokens, 8 candidates, 1 question, p50 / p95 & \textbf{50 / 50} & 55 / 59 & 29 / 29 & 30 / 30 \\
W2: 3.2k-token state, 8 questions, one batch & 550 & 709 & 137 & 167 \\
W2: parent question first, 7 reuse the prefix cache & 617 & 736 & 179 & 201 \\
W3: 8,192 tokens, 1 question & 417 & 698 & 120 & 148 \\
W3: 16,320 tokens, 1 question & 606 & 766 & 176 & 230 \\
W4: best throughput, input tokens/s & \textbf{16,282} & 12,524 & 54,812 & 41,794 \\
\bottomrule
\end{tabular}
\end{adjustbox}
\end{table}

\paragraph{The engine lowers short-request latency.}
W1 takes 50\,ms at the median in FP8 and 55\,ms in BF16 (Table~\ref{tab:engine}), against
97\,ms for the transformers path on one B300 (Section~\ref{sec:serving}). Throughput is
similar: 16,282 input tokens per second for vLLM FP8 on the H100 against 16,072 for
transformers on the B300. The GPUs differ, so this is not a controlled engine comparison.
One tail value stands out: BF16 at 16K tokens has a p95 of 1,410\,ms against a p50 of 766\,ms.

\paragraph{The engine's prefix cache did not help W2.}
vLLM caches the hybrid layers' recurrent state in its \texttt{align} mode. Running the parent
question first and letting the other seven reuse its cached prefix was slower than submitting
all eight as one batch (617 against 550\,ms in FP8). Our own state branching on the
transformers path (Section~\ref{sec:branching}) reduced W2 from 1,633 to 379\,ms on the B300;
it has not been ported to vLLM.

\paragraph{A 9B backbone is much faster.}
Qwen3.5-9B, measured for speed only and not trained for Bobcat, reaches 54,812 input tokens
per second in FP8 and answers W1 in 29\,ms. Its decision quality after training is unknown;
zero-shot it scored 78.4\% on development (Section~\ref{sec:leaderboard}).

\paragraph{Remaining engine headroom.}
Porting our hybrid state branching to vLLM would, by our estimate, make shared-context
workloads three to four times faster; CUDA graphs, an FP8 KV cache and larger-batch
scheduling also remain untried. None of these is measured here.

\subsection{What was not measured}

HTTP handling, the network and the authenticating edge are excluded from these timings. On the
internal deployment's L40S, a W1-shaped request (588 input tokens, eight candidates) took
154\,ms at the median over local HTTP, and a three-question request from one client through
the edge over the internet took 294\,ms at the median over ten requests; latency under load
has not been measured. W4 measures engine throughput, not latency under a queue of
concurrent users. The 9B numbers are speed only.
